Sequential LoRA adaptation (SeqLoRA) suffers from catastrophic forgetting because a single
shared adapter drifts away from old-task optima with no structural protection.
Two families of remedies exist---sharpness-aware optimization (flatness) and
Fisher-weighted regularization (stability)---but their geometric interaction in the LoRA
continual learning setting has not been studied.

We ask: do these two objectives \emph{conflict} or \emph{complement} each other in gradient space?
We show that the flatness gradient component (adapter-subspace first-order sharpness, AS$^1$)
and the past-task Fisher gradient are \textbf{near-orthogonal}
($\cosff \approx -0.03$ to $-0.001$, verified across six fine-grained recognition datasets).
This orthogonality arises from spectral separation between the current-task Hessian
and the past-task Fisher in the effective adapter update space $\DeltaW = BA$,
amplified by the Jacobian of the LoRA reparametrization.

Building on this finding, we propose \textbf{Flat-Stable LoRA (FS-LoRA)},
which (1)~computes AS$^1$ flatness in adapter factor coordinates over the current-task loss only,
(2)~applies trace-normalized Fisher-EWC as an independent gradient component in $\DeltaW$-space,
and (3)~combines them via an orthogonality-preserved decoupled update that prevents the
two objectives from coupling in a shared perturbation closure.
Experiments on six fine-grained recognition datasets show improvements on datasets where
the Fisher penalty is active (Flowers $+5.5$pp, CUB200 $+1.0$pp, ImageNet-R $+0.2$pp),
with mixed results on Cars196 and Aircraft; all results are single-seed.
On standard PECL benchmarks, LR-tuned results show FS-LoRA within $1$pp of EWC-LoRA
on ImageNet-A ($58.92$ vs.~$59.89$), while CIFAR-100, ImageNet-R, and DomainNet
runs are still in progress.
